\subsection{Defining the problem}

Observations $Y_1, \ldots, Y_n$ arise in time order. 

\begin{enumerate}
    \item Starting from
    $$f_{Y_1, \ldots, Y_n}(y_1, \ldots, y_n) \; = \; f_{Y_n | Y_1, \ldots, Y_{n-1}}(y_n | y_1, \ldots, y_{n-1})f_{Y_1, \ldots, Y_{n-1}}(y_1, \ldots, y_{n-1}),$$
    establish the \textit{prediction decomposition}
    $$f_{Y_1, \ldots, Y_n}(y_1, \ldots, y_n) = f_{Y_1}(y_1) \prod_{j=2}^{n}f_{Y_j | Y_1, \ldots, Y_{j-1}}(y_j | y_1, \ldots, y_{j-1}).$$
    \item A stationary first-order Gaussian autoregressive process satisfies
    $$Y_j | Y_1 = y_1, \ldots, Y_j | Y_{j-1} = y_{j-1} \sim \mathcal{N}(\mu + \alpha(y_{j-1}-\mu), \sigma^2) \; \text{for } j = 1, \ldots, n$$
    with $|\alpha| < 1$, $\mu \in \mathbb{R}$, and $\sigma^2 > 0$. Find the log likelihood for data $y_0, \ldots, y_n$ for this model if the initial value $y_0$ is treated
    \begin{enumerate}
        \item fixed as known constant,
        \item coming from the stationary distribution $\mathcal{N}(\mu, \frac{\sigma^2}{1-\alpha^2})$
    \end{enumerate}
    \item Give a minimal sufficient statistic for $\theta = (\mu, \sigma^2, \alpha)$ in (2). Does the model with parameter vector $\theta$ form
    an exponential family?
\end{enumerate}

\subsection{Solution}

We first show the prediction decomposition formula starting from the given identity on (1).
\subsubsection{Prediction decomposition derivation}

$$f_{Y_1, \ldots, Y_n}(y_1, \ldots, y_n) \; = \; f_{Y_n | Y_1, \ldots, Y_{n-1}}(y_n | y_1, \ldots, y_{n-1})f_{Y_1, \ldots, Y_{n-1}}(y_1, \ldots, y_{n-1})$$
$$= f_{Y_n | Y_1, \ldots, Y_{n-1}}(y_n | y_1, \ldots, y_{n-1}) f_{Y_{n-1} | Y_1, \ldots, Y_{n-2}}(y_{n-1} | y_1, \ldots, y_{n-2}) f_{Y_1, \ldots, Y_{n-2}}(y_1, \ldots, y_{n-2})$$
$$= f_{Y_n | Y_1, \ldots, Y_{n-1}}(y_n | y_1, \ldots, y_{n-1}) f_{Y_{n-1} | Y_1, \ldots, Y_{n-2}}(y_{n-1} | y_1, \ldots, y_{n-2}) \cdots f_{Y_2 | Y_1}(y_2 | y_1) f_{Y_1}(y_1)$$
$$= f_{Y_1}(y_1) \prod_{j=2}^{n}f_{Y_j | Y_1, \ldots, Y_{j-1}}(y_j | y_1, \ldots, y_{j-1}).$$ 

This shows that the joint density equals the product of the onestep ahead predictive densities. In other words, the joiunt behaviour of the whole sequience is built by multiplying all the
``predictive'' pieces.

\subsubsection{Likelihood functions for AR(1) process}

Now, suppose we observe data $y_1, \ldots, y_n$ from the autoregressive process
\[
Y_j = \mu + \alpha(Y_{j-1} - \mu) + \varepsilon_j, 
\qquad \varepsilon_j \sim \mathcal{N}(0, \sigma^2),
\]
with $|\alpha| < 1$. Note that the \textbf{gaussian} noises $\varepsilon_j$'s are i.i.d. This model means each $Y_j$ depends linearly on the previous value $Y_{j-1}$ plus some gaussian noise.
Hence, this implies that 
$$Y_j | Y_{j-1} \sim \mathcal{N}(\mu + \alpha(y_{j-1} - \mu), \sigma^2).$$ 

\paragraph{Case (i): Dealing with $y_0$ as a known constant.}
If the initial value $y_0$ is treated as a fixed, known constant, then we may write the joint pdf of the realizations ($y_1, \ldots, y_n$) of the observations as:

$$f(y_1, \ldots, y_n | \mu, \alpha, \sigma^2) = \prod_{j=1}^{n} f(y_j | y_{j-1};\mu, \alpha, \sigma^2)$$

And so, we might write the conditional log-likelihood as
\[
\ell(\mu,\alpha,\sigma^2)
  = -\frac{n}{2}\log(2\pi\sigma^2)
    -\frac{1}{2\sigma^2}\sum_{j=1}^n
       [y_j-\mu-\alpha(y_{j-1}-\mu)]^2.
\]
Here $y_0$ enters only in the first term (i.e. the term $y_1 \mid y_0$),
but no density is assigned to $y_0$ by itself.

\paragraph{Case (ii): Dealing witht the devil; $y_0$ is random, from stationary distribution.}
If we assume the process is stationary, then
\[
Y_0 \sim \mathcal{N}\!\left(\mu,\,\frac{\sigma^2}{1-\alpha^2}\right),
\]
and so, weget that the marginal (stationary) distribution of the first observation is
\[
Y_1 \sim \mathcal{N}\!\left(\mu,\,\frac{\sigma^2}{1-\alpha^2}\right).
\]
Integrating out the unobserved $Y_0$ leads to the \emph{exact likelihood}
\[
L_e(\mu,\alpha,\sigma^2 \,;\, y_1,\ldots,y_n)
  = f\!\left(y_1; \mu, \tfrac{\sigma^2}{1-\alpha^2}\right)
    \prod_{j=2}^n f(y_j \mid y_{j-1};\, \mu,\alpha,\sigma^2),
\]
and the corresponding log-likelihood is
$$
\ell_e(\mu,\alpha,\sigma^2)
  = -\frac{1}{2}\log\!\frac{2\pi\sigma^2}{1-\alpha^2}
    -\frac{(y_1-\mu)^2(1-\alpha^2)}{2\sigma^2}
    -\frac{n-1}{2}\log(2\pi\sigma^2)
    -\frac{1}{2\sigma^2}\sum_{j=2}^n
       [y_j-\mu-\alpha(y_{j-1}-\mu)]^2.
$$

\paragraph{Interpretation (for me to remember $\ldots$):}
\begin{itemize}
  \item The \textbf{conditional likelihood} treats $y_0$ as a fixed starting point and models only the transitions $Y_j \mid Y_{j-1}$ for $j=1,\dots,n$.
  \item The \textbf{exact likelihood} assumes $Y_0$ (and hence $Y_1$) come from the stationary distribution and therefore includes their marginal densities.
  \item Both forms give consistent estimates of $(\mu,\alpha,\sigma^2)$ as $n\to\infty$, but they differ slightly in small samples.
\end{itemize}

\subsubsection{Minimal Sufficient Statistic}
Note that likelihood depends on the data only through (if expanded and rearranged, the term in the summation that is)
$$\hat{T}(y)=\left(\sum_{j=1}^n y_j, \sum_{j=1}^n y_{j-1}, \sum_{j=1}^n y_j y_{j-1}, \sum_{j=1}^n y_j^2, \sum_{j=1}^n y_{j-1}^2\right),$$

\begin{remark}
    Note that the term $\sum_{j=1}^n y_{j-1}^2$ can be ``omitted" (more like rewritten) since it can be computed from $\sum_{j=1}^n y_j^2$ and the known initial value $y_0$ (i.e. it is just 
    lagged term of the $y_j$) as follows:
    $$\sum_{j=1}^n y_{j-1}^2 = \sum_{j=0}^{n-1} y_j^2 = \sum_{j=1}^{n} y_j^2 + y_0^2 - y_n^2.$$
\end{remark}

following the remark assbove, we define the following statistic:
\vspace{0.2cm}
$$T(y)=\left(\sum_{j=1}^n y_j, \sum_{j=1}^n y_{j-1}, \sum_{j=1}^n y_j y_{j-1}, \sum_{j=1}^n y_j^2\right),$$
\vspace{0.2cm}

which is then minimal sufficient for $\theta=(\mu,\sigma^2,\alpha)$.

\subsubsection*{Exponential Family}

\begin{theorem}[Conditional AR(1) is an exponential family]
Fix $y_0\in\mathbb R$. Let $(Y_j)_{j=1}^n$ satisfy the Gaussian AR(1)
\[
Y_j \mid Y_{j-1}=y_{j-1}\ \sim\ \mathcal N\!\big(\mu+\alpha(y_{j-1}-\mu),\,\sigma^2\big),
\quad j=1,\dots,n,
\]
with parameter $\theta=(\mu,\alpha,\sigma^2)$ and $|\alpha|<1$. Then the conditional joint density
$f_\theta(y_1,\dots,y_n\mid y_0)=\prod_{j=1}^n f_\theta(y_j\mid y_{j-1})$ belongs to a (regular) exponential family.
\end{theorem}

\begin{proof}
Write the conditional density
\[
f_\theta(y_j\mid y_{j-1})
=(2\pi\sigma^2)^{-1/2}
\exp\!\left\{-\frac{1}{2\sigma^2}\big(y_j-\mu-\alpha(y_{j-1}-\mu)\big)^2\right\}.
\]
Hence
\[
\log f_\theta(y_{1:n}\mid y_0)
=-\frac{n}{2}\log(2\pi\sigma^2)
-\frac{1}{2\sigma^2}\sum_{j=1}^n\big(y_j-\alpha y_{j-1}-(1-\alpha)\mu\big)^2.
\]
Expanding the square and collecting terms that depend on the data gioves

\[
\log f_\theta(y_{1:n}\mid y_0)
= C(\theta) + \eta_1(\theta)\sum_{j=1}^n y_j^2
+\eta_2(\theta)\sum_{j=1}^n y_{j-1}^2
+\eta_3(\theta)\sum_{j=1}^n y_jy_{j-1}
+\eta_4(\theta)\sum_{j=1}^n y_j
+\eta_5(\theta)\sum_{j=1}^n y_{j-1},
\]

where the \emph{natural parameters} $(\eta_k(\theta))$ are smooth functions of $(\mu,\alpha,\sigma^2)$
(e.g.\ $\eta_1(\theta)=-\tfrac{1}{2\sigma^2}$, $\eta_3(\theta)=\tfrac{\alpha}{\sigma^2}$, etc.), and
$C(\theta)$ \emph{absorbs} all terms free of $y_{1:n} = (y_1, \ldots, y_n)$.
So
\[
f_\theta(y_{1:n}\mid y_0)
= h(y_{1:n}\mid y_0)\,
\exp\!\Big\{ \sum_{k=1}^5 \eta_k(\theta)\, T_k(y_{1:n}) - A(\theta)\Big\},
\]
where $h(y_{1:n}\mid y_0)\equiv 1$ (I believe this function is sometimes called \emph{carrier}), and sufficient statistics
\[
T_1=\sum y_j^2,\quad
T_2=\sum y_{j-1}^2,\quad
T_3=\sum y_jy_{j-1},\quad
T_4=\sum y_j,\quad
T_5=\sum y_{j-1}.
\]

This is exactly the canonical (normal, `bread and butter') exponential-family form in $\mathbb R^5$, hence the claim. (Note that $-C(\theta) = A(\theta)$)
\end{proof}

\begin{remark}[On the role of $Y_0$]
The statement is \emph{conditional on} $y_0$. If one assumes a stationary start
$Y_0\sim\mathcal N(\mu,\sigma^2/(1-\alpha^2))$, then the \emph{exact} joint density
$f_\theta(y_0,\dots,y_n)$ is also an exponential family after adding the marginal term for $y_0$,
with the same type of quadratic/linear sufficient statistics augmented by $y_0$-terms.
If $Y_0$ has unspecified law, only the conditional likelihood admits the exponential-family factorization.
\end{remark}
